% Options for packages loaded elsewhere
\PassOptionsToPackage{unicode}{hyperref}
\PassOptionsToPackage{hyphens}{url}
\PassOptionsToPackage{dvipsnames,svgnames,x11names}{xcolor}
\PassOptionsToPackage{space}{xeCJK}
\documentclass[
]{article}
\usepackage{xcolor}
\usepackage[margin=2.5cm]{geometry}
\usepackage{amsmath,amssymb}
\setcounter{secnumdepth}{-\maxdimen} % remove section numbering
\usepackage{iftex}
\ifPDFTeX
  \usepackage[T1]{fontenc}
  \usepackage[utf8]{inputenc}
  \usepackage{textcomp} % provide euro and other symbols
\else % if luatex or xetex
  \usepackage{unicode-math} % this also loads fontspec
  \defaultfontfeatures{Scale=MatchLowercase}
  \defaultfontfeatures[\rmfamily]{Ligatures=TeX,Scale=1}
\fi
\usepackage{lmodern}
\ifPDFTeX\else
  % xetex/luatex font selection
  \setmainfont[]{DejaVu Sans}
  \setmonofont[]{DejaVu Sans Mono}
  \ifXeTeX
    \usepackage{xeCJK}
    \setCJKmainfont[]{Noto Sans CJK SC}
  \fi
  \ifLuaTeX
    \usepackage[]{luatexja-fontspec}
    \setmainjfont[]{Noto Sans CJK SC}
  \fi
\fi
% Use upquote if available, for straight quotes in verbatim environments
\IfFileExists{upquote.sty}{\usepackage{upquote}}{}
\IfFileExists{microtype.sty}{% use microtype if available
  \usepackage[]{microtype}
  \UseMicrotypeSet[protrusion]{basicmath} % disable protrusion for tt fonts
}{}
\makeatletter
\@ifundefined{KOMAClassName}{% if non-KOMA class
  \IfFileExists{parskip.sty}{%
    \usepackage{parskip}
  }{% else
    \setlength{\parindent}{0pt}
    \setlength{\parskip}{6pt plus 2pt minus 1pt}}
}{% if KOMA class
  \KOMAoptions{parskip=half}}
\makeatother
\usepackage{longtable,booktabs,array}
\usepackage{caption}
\captionsetup[table]{skip=6pt}
\newcounter{none} % for unnumbered tables
\usepackage{calc} % for calculating minipage widths
% Correct order of tables after \paragraph or \subparagraph
\usepackage{etoolbox}
\makeatletter
\patchcmd\longtable{\par}{\if@noskipsec\mbox{}\fi\par}{}{}
\makeatother
% Allow footnotes in longtable head/foot
\IfFileExists{footnotehyper.sty}{\usepackage{footnotehyper}}{\usepackage{footnote}}
\makesavenoteenv{longtable}
\setlength{\emergencystretch}{3em} % prevent overfull lines
\providecommand{\tightlist}{%
  \setlength{\itemsep}{0pt}\setlength{\parskip}{0pt}}
\usepackage{bookmark}
\IfFileExists{xurl.sty}{\usepackage{xurl}}{} % add URL line breaks if available
\urlstyle{same}
% fallback for those not using the hyperref driver hyperxmp:
\makeatletter
\@ifundefined{xmpquote}{\newcommand{\xmpquote}[1]{#1}}{}
\makeatother
\hypersetup{
  colorlinks=true,
  linkcolor={Maroon},
  filecolor={Maroon},
  citecolor={Blue},
  urlcolor={Blue},
  pdfcreator={LaTeX via pandoc}}

\author{}
\date{}

\begin{document}

\section{筑基篇课后习题 V：pat
重新观测杨-米尔斯存在性与质量间隙}\label{ux7b51ux57faux7bc7ux8bfeux540eux4e60ux9898-vpat-ux91cdux65b0ux89c2ux6d4bux6768-ux7c73ux5c14ux65afux5b58ux5728ux6027ux4e0eux8d28ux91cfux95f4ux9699}

\begin{quote}
\textbf{声明 (Declaration)}: 本文\textbf{不代表任何数学结论}
(non-mathematical-claim), 仅为基于 Lean 4
形式化的一种\textbf{实验观测报告} (experimental observation report)。

{\def\LTcaptype{none} % do not increment counter
\begin{longtable}[]{@{}
  >{\raggedright\arraybackslash}p{(\linewidth - 6\tabcolsep) * \real{0.2500}}
  >{\raggedright\arraybackslash}p{(\linewidth - 6\tabcolsep) * \real{0.2500}}
  >{\raggedright\arraybackslash}p{(\linewidth - 6\tabcolsep) * \real{0.2500}}
  >{\raggedright\arraybackslash}p{(\linewidth - 6\tabcolsep) * \real{0.2500}}@{}}
\toprule\noalign{}
\begin{minipage}[b]{\linewidth}\raggedright
习题
\end{minipage} & \begin{minipage}[b]{\linewidth}\raggedright
主题
\end{minipage} & \begin{minipage}[b]{\linewidth}\raggedright
Lean 形式化
\end{minipage} & \begin{minipage}[b]{\linewidth}\raggedright
DOI
\end{minipage} \\
\midrule\noalign{}
\endhead
\bottomrule\noalign{}
\endlastfoot
exercise\_i\_3 & pat 重新观测杨-米尔斯（质量间隙） & PatYangMills.lean &
10.5281/zenodo.21916843 \\
\end{longtable}
}
\end{quote}

\textbf{Exercise V for the Foundation-Building Chapter: Yang--Mills
Existence and Mass Gap Revisited from the PAT Perspective}

\begin{quote}
习题定位：筑基篇课后习题系列第五题。用筑基篇的相位/互锁/折叠结构（发散-周期特征分解、相位-数值互锁、单位根相位圆、折叠类
0）重新观测杨-米尔斯（Yang--Mills）存在性与质量间隙问题（Clay
千禧年问题之一，Jaffe--Witten
陈述）。全部新定理只锚筑基篇定理与框架（R047/R085/R110/R139/R141），不用外部引理。诚实边界延续：\textbf{不证明质量间隙存在}（该问题的存在性陈述本身就是千禧年开放问题）。
\end{quote}

\emph{2026-08-13 · Internal research exercise · Lean 4 / mathlib v4.32.2
· 7 new theorems PROVED no sorry · 算器神魂论筑基篇课后习题 V }

\begin{center}\rule{0.5\linewidth}{0.5pt}\end{center}

\subsection{习题陈述}\label{ux4e60ux9898ux9648ux8ff0}

\begin{quote}
站在 pat
视角下重新观测杨-米尔斯存在性与质量间隙：用筑基篇已证定理（R047/R085/R110/R139/R141）把经典
YM 的无质量性与质量间隙的谱结构重述为 pat 几何。唯一论点：\textbf{经典
YM 无质量 = 纯相位（无数值锁定）；质量 =
相位-数值互锁的数值侧（R139）；质量间隙 = 折叠类
0（真空）与第一非平凡相位层的间隙------最短非平凡相位圆。}
\end{quote}

\begin{center}\rule{0.5\linewidth}{0.5pt}\end{center}

\subsection{零、总论：为什么质量间隙在 pat
视角下是相位圆与折叠类的故事}\label{ux96f6ux603bux8bbaux4e3aux4ec0ux4e48ux8d28ux91cfux95f4ux9699ux5728-pat-ux89c6ux89d2ux4e0bux662fux76f8ux4f4dux5706ux4e0eux6298ux53e0ux7c7bux7684ux6545ux4e8b}

筑基篇 R047 钉死发散/周期的结构：实轴（发散）= 共轭对称 S
的固定分量，虚轴 J（周期）= S 的反射分量，二者共享基点 0、方向正交。R139
钉死相位-数值互锁：声明必须是两组对称性（相位/方向对 +
数值/距离对）。R085 钉死折叠类：0 = 折叠中心。

杨-米尔斯质量间隙在 pat 视角下是这三者的合成：

\begin{itemize}
\tightlist
\item
  \textbf{经典 YM 无质量} = 纯相位：经典理论尺度不变（无质量参数），相位
  e\^{}\{-imt\} 在 m = 0 时冻结在
  1------相位不演进，无数值变化（R047：发散轴上无周期）。
\item
  \textbf{质量 = 相位周期倒数}：质量 m 的相位以 T = 2π/m
  为周期，绕相位圆一圈回到基点 1（R141 单位根 n 槽环的连续版；T·m = 2π
  互锁）。
\item
  \textbf{质量间隙} = 折叠类 0（真空）与第一非平凡相位层的间隙：真空 =
  质量 0 = 折叠类
  0（R085）；最小正质量的相位周期严格为正------第一非平凡层与折叠类的间隙非零。
\end{itemize}

\begin{center}\rule{0.5\linewidth}{0.5pt}\end{center}

\subsection{一、经典 YM =
纯相位（无数值锁定）（R047）}\label{ux4e00ux7ecfux5178-ym-ux7eafux76f8ux4f4dux65e0ux6570ux503cux9501ux5b9ar047}

\textbf{★核心：经典 YM 无质量（尺度不变）------相位 e\^{}\{-imt\} 在 m =
0 时冻结在 1，相位不演进，无数值变化。这是 R047
发散轴语义的物理实例：无质量 = 发散轴上无周期。}

\subsubsection{论证}\label{ux8bbaux8bc1}

\begin{enumerate}
\def\labelenumi{\arabic{enumi}.}
\tightlist
\item
  \textbf{经典 YM 尺度不变}（Jaffe--Witten 陈述的已知部分）：经典
  Yang--Mills
  拉氏量没有质量参数，理论无固有能量/时间尺度------规范玻色子无质量，传播子
  1/k² 沿实轴（发散）方向。
\item
  \textbf{无质量 = 相位冻结}：质量 m 的平面波相位 e\^{}\{-imt\}（ω =
  m，ħ = c = 1）；m = 0 时相位恒为
  1------相位不转（不演进），无数值变化。这正是''无质量 =
  无固有周期''的相位表达。
\item
  \textbf{R047 锚定}：周期结构需要周期轴 J（发散轴上 -1 无平方根，R047
  divergence\_axis\_no\_sqrt）；无质量粒子在发散轴上，无周期------纯相位、纯发散。
\item
  \textbf{诚实边界}：经典 YM
  无质量是已知事实；量子理论的质量生成（dimensional transmutation，格点
  QCD 数值证据）不在本题------本题形式化的是相位结构本身。
\end{enumerate}

\subsubsection{形式化（PatYangMills.lean，1
定理）}\label{ux5f62ux5f0fux5316patyangmills.lean1-ux5b9aux7406}

\begin{itemize}
\tightlist
\item
  \texttt{massless\_phase\_frozen}：\texttt{∀\ t,\ Complex.exp\ (-(0·t)·i)\ =\ 1}------无质量相位冻结（相位不演进）。
\end{itemize}

\textbf{验证}：0 sorry。simp 直接闭合（0 吸收律，heap\_verify Z/7
枚举验证代数方向）。

\begin{center}\rule{0.5\linewidth}{0.5pt}\end{center}

\subsection{二、质量 = 相位周期的倒数（R141
连续版）}\label{ux4e8cux8d28ux91cf-ux76f8ux4f4dux5468ux671fux7684ux5012ux6570r141-ux8fdeux7eedux7248}

\textbf{★核心：质量 m 的相位 e\^{}\{-imt\} 以 T = 2π/m
为周期（绕相位圆一圈回到基点 1）；质量与周期互锁 T·m = 2π。这是 R141
单位根 n 槽环的连续版：离散版 n 个相位均匀分布绕圆，连续版质量相位以
2π/m 为最小圆。}

\subsubsection{论证}\label{ux8bbaux8bc1-1}

\begin{enumerate}
\def\labelenumi{\arabic{enumi}.}
\tightlist
\item
  \textbf{相位往返}：e\^{}\{-im(t+T)\} = e\^{}\{-imt\} ⟺ mT = 2π ⟺ T =
  2π/m（m ≠ 0）。绕相位圆一圈（相位差 2π）回到基点
  1------CompactToolkit.exp\_two\_pi\_I\_eq\_one：exp(2πi) = 1。
\item
  \textbf{R141 锚定}：pat n 蜷曲到圆 + 单位根量化（误差 ≤
  π/n）------离散 n 槽环；质量相位的连续版 = 相位圆，质量周期 T = 2π/m
  是绕圆一周的时间。
\item
  \textbf{质量在周期轴上}：质量 m 的相位分量沿周期轴 J（R047：J
  是共轭对称 S 的反射方向）；m 倍保持方向------周期轴 = S 的 -1
  特征空间，与发散轴共享基点 0、正交。
\item
  \textbf{正质量 ⟹ 正周期}：m \textgreater{} 0 ⟹ 2π/m \textgreater{}
  0------最小正质量对应最短非平凡相位圆（见第四节）。
\end{enumerate}

\subsubsection{形式化（PatYangMills.lean，3
定理）}\label{ux5f62ux5f0fux5316patyangmills.lean3-ux5b9aux7406}

\begin{itemize}
\tightlist
\item
  \texttt{mass\_phase\_period}：\texttt{Complex.exp\ (-(m·(t\ +\ 2π/m))·i)\ =\ Complex.exp\ (-(m·t)·i)}（m
  ≠ 0）------相位往返，周期 T = 2π/m。
\item
  \texttt{mass\_period\_axis}：\texttt{conj\ ⟨0,\ m⟩\ =\ -⟨0,\ m⟩}------质量分量在周期轴
  J 上（S 反射方向）。
\item
  \texttt{positive\_mass\_positive\_period}：\texttt{0\ \textless{}\ m\ →\ 0\ \textless{}\ 2π/m}------正质量
  = 正相位周期。
\end{itemize}

\textbf{验证}：0 sorry。\texttt{mass\_phase\_period} 用 field\_simp（m ≠
0）+ ring + exp\_sub/exp\_neg +
\texttt{CompactToolkit.exp\_two\_pi\_I\_eq\_one}；教训：\texttt{-(2π)·i}
不是 exp(-x) 模式，需先 ring 化为 \texttt{-((2π)·i)} 再 rw
{[}Complex.exp\_neg{]}。

\begin{center}\rule{0.5\linewidth}{0.5pt}\end{center}

\subsection{三、质量对（m, 1/m）= log
镜像互锁（R139/R110）}\label{ux4e09ux8d28ux91cfux5bf9m-1m-log-ux955cux50cfux4e92ux9501r139r110}

\textbf{★核心：Compton 波长 λ = 1/m------质量与长度互为倒数对，log
镜像对称（log(1/m) = -log m）。质量 =
相位-数值互锁（R139）的数值侧：相位锁定（e\^{}\{-imt\}
的圆）给出数值锁定（m 的倒数 λ）。}

\subsubsection{论证}\label{ux8bbaux8bc1-2}

\begin{enumerate}
\def\labelenumi{\arabic{enumi}.}
\tightlist
\item
  \textbf{R139 互锁}：声明 = 两组对称性（相位/方向对 +
  数值/距离对）；相位锁定的数值侧 = 距离（波长）。成对向量 = 矩阵；互锁
  ⟺ 矩阵非奇异（相位 ↔ 数值双向可解）。
\item
  \textbf{Compton 波长}：λ = 1/m------质量与长度互为倒数对（自然单位 ħ =
  c = 1）。
\item
  \textbf{R110 log 镜像}：数值对 (r, 1/r) 的对称性是 log
  镜像对称（log(1/a) = -log a）；质量对 (m, 1/m) 是 R139
  magnitude\_pair\_log\_mirror 的直接实例。
\item
  \textbf{pat 视角}：质量的相位侧（周期 T = 2π/m）与数值侧（波长 λ =
  1/m）经 2π
  与互锁互锁------质量既锁定相位圆（周期）又锁定距离圆（Compton
  波长），两组对称性一体的数值侧。
\end{enumerate}

\subsubsection{形式化（PatYangMills.lean，1
定理）}\label{ux5f62ux5f0fux5316patyangmills.lean1-ux5b9aux7406-1}

\begin{itemize}
\tightlist
\item
  \texttt{mass\_compton\_log\_mirror}：\texttt{0\ \textless{}\ m\ →\ log(1/m)\ =\ -log\ m}------质量对
  log 镜像（锚 \texttt{MutualLocking.magnitude\_pair\_log\_mirror}）。
\end{itemize}

\textbf{验证}：0 sorry。直接锚 R139 已证定理（rw {[}one\_div,
Real.log\_inv{]} 已在 MutualLocking 内完成）。

\begin{center}\rule{0.5\linewidth}{0.5pt}\end{center}

\subsection{四、质量间隙 = 折叠类
0（真空）与第一非平凡相位层（R085）}\label{ux56dbux8d28ux91cfux95f4ux9699-ux6298ux53e0ux7c7b-0ux771fux7a7aux4e0eux7b2cux4e00ux975eux5e73ux51e1ux76f8ux4f4dux5c42r085}

\textbf{★核心：真空 = 质量 0 = 折叠类 0（R085：mirror fixes 0，0
是折叠中心）；质量间隙 Δ \textgreater{} 0 的结构对应 = 折叠类 0
与第一非平凡相位层（最小正质量）的间隙------最小正质量的相位周期严格为正。诚实边界：这是几何对应，非存在性证明。}

\subsubsection{论证}\label{ux8bbaux8bc1-3}

\begin{enumerate}
\def\labelenumi{\arabic{enumi}.}
\tightlist
\item
  \textbf{真空 = 折叠类 0}（R085）：mirror fixes 0------0
  是折叠中心；真空（质量 0，无粒子）在折叠类 0。质量间隙的物理意义 =
  真空与最低激发态的能量差严格为正（Jaffe--Witten 陈述：最低激发态能量
  \textgreater{} 0）。
\item
  \textbf{第一非平凡相位层}：最小正质量 m\_min 的相位周期 T\_min =
  2π/m\_min
  严格为正（positive\_mass\_positive\_period）------第一非平凡层与折叠类
  0 的间隙非零（0 \textless{} T\_min）。
\item
  \textbf{实数稠密性陷阱（已避开）}：``最小非零距离''在 ℝ
  上不存在（无最小正实数）------质量间隙不能形式化为 min
  非零距；正确形式化 = 最小正质量 ⟹ 正相位周期（结构对应）。
\item
  \textbf{诚实边界}：质量间隙的存在性（量子 YM
  的最低激发态能量严格为正）是千禧年开放问题------本题交付的是其 pat
  几何对应（折叠类 0 与第一非平凡相位层的间隙），非存在性证明。格点 QCD
  数值证据（Wilson 1974）支持质量间隙，但严格数学证明（Jaffe--Witten
  问题）仍开放。
\end{enumerate}

\subsubsection{形式化（PatYangMills.lean，1
定理）}\label{ux5f62ux5f0fux5316patyangmills.lean1-ux5b9aux7406-2}

\begin{itemize}
\tightlist
\item
  \texttt{vacuum\_fold\_class}：\texttt{-(0)\ =\ 0}------真空（质量 0）=
  折叠类 0（锚 \texttt{MirrorFoldZero.mirror\_fixes\_zero}）。
\end{itemize}

\textbf{验证}：0 sorry。与
positive\_mass\_positive\_period（第二节）共同给出间隙的结构对应。

\begin{center}\rule{0.5\linewidth}{0.5pt}\end{center}

\subsection{五、全景：经典无质量 ∧ 质量 = 周期倒数 ∧ 质量对 log 镜像 ∧
真空折叠类（组合定理）}\label{ux4e94ux5168ux666fux7ecfux5178ux65e0ux8d28ux91cf-ux8d28ux91cf-ux5468ux671fux5012ux6570-ux8d28ux91cfux5bf9-log-ux955cux50cf-ux771fux7a7aux6298ux53e0ux7c7bux7ec4ux5408ux5b9aux7406}

\textbf{★核心：无质量相位冻结（经典 YM = 纯相位）∧ 质量 =
相位周期倒数（T·m = 2π）∧ 质量对 log 镜像（Compton）∧ 真空 = 折叠类 0 ∧
正质量正周期------质量 = 相位-数值互锁的数值侧（R139），质量间隙 =
折叠类 0 与第一非平凡相位层的间隙。}

\subsubsection{形式化（PatYangMills.lean，1
定理）}\label{ux5f62ux5f0fux5316patyangmills.lean1-ux5b9aux7406-3}

\begin{itemize}
\tightlist
\item
  \texttt{ym\_pat\_perspective}：全景------五层组合（逐项锚一至四节定理）。
\end{itemize}

\textbf{验证}：0 sorry。合取项分别锚 massless\_phase\_frozen /
mass\_phase\_period / mass\_compton\_log\_mirror / vacuum\_fold\_class /
positive\_mass\_positive\_period。

\begin{center}\rule{0.5\linewidth}{0.5pt}\end{center}

\subsection{习题解答总结}\label{ux4e60ux9898ux89e3ux7b54ux603bux7ed3}

{\def\LTcaptype{none} % do not increment counter
\begin{longtable}[]{@{}
  >{\raggedright\arraybackslash}p{(\linewidth - 6\tabcolsep) * \real{0.2500}}
  >{\raggedright\arraybackslash}p{(\linewidth - 6\tabcolsep) * \real{0.2500}}
  >{\raggedright\arraybackslash}p{(\linewidth - 6\tabcolsep) * \real{0.2500}}
  >{\raggedright\arraybackslash}p{(\linewidth - 6\tabcolsep) * \real{0.2500}}@{}}
\toprule\noalign{}
\begin{minipage}[b]{\linewidth}\raggedright
论点
\end{minipage} & \begin{minipage}[b]{\linewidth}\raggedright
内容
\end{minipage} & \begin{minipage}[b]{\linewidth}\raggedright
锚到
\end{minipage} & \begin{minipage}[b]{\linewidth}\raggedright
定理
\end{minipage} \\
\midrule\noalign{}
\endhead
\bottomrule\noalign{}
\endlastfoot
经典 YM = 纯相位 & 无质量相位冻结（相位不演进） & R047 + Jaffe--Witten
已知部分 & massless\_phase\_frozen \\
质量 = 周期倒数 & T = 2π/m 相位往返；质量在周期轴 J 上 & R141 连续版 +
R047 & mass\_phase\_period / mass\_period\_axis \\
质量对 log 镜像 & Compton λ = 1/m（log 镜像） & R139/R110 &
mass\_compton\_log\_mirror \\
真空 = 折叠类 0 & 质量 0 = 折叠中心；正质量正周期 & R085 &
vacuum\_fold\_class / positive\_mass\_positive\_period \\
★全景 & 质量 = 互锁数值侧；间隙 = 折叠类 0 与第一非平凡层 &
R047/R085/R139/R141 & ym\_pat\_perspective \\
\end{longtable}
}

\textbf{习题的回答（一句话）}：\textbf{pat
视角下杨-米尔斯质量间隙是相位圆与折叠类的故事------经典 YM 无质量 =
纯相位（相位冻结，发散轴）；质量 = 相位周期倒数（绕相位圆一圈 = T·m =
2π，周期轴 J）；质量间隙 = 折叠类
0（真空）与第一非平凡相位层的间隙（最小正质量 =
最短非平凡相位圆）。这是结构侧写，不是千禧年问题证明。}

\textbf{诚实边界}（延续全部习题）：质量间隙存在性（量子 YM
最低激发态能量严格为正，Jaffe--Witten）\textbf{未证明}------本题交付 pat
结构侧写（折叠类 0
与第一非平凡相位层的间隙、相位-数值互锁的数值侧），全部 7
个新定理是结构事实的 Lean 验收，非质量间隙存在性断言。经典 YM
无质量（尺度不变）是已知事实。

\begin{center}\rule{0.5\linewidth}{0.5pt}\end{center}

\subsection{六、深入（R162）：相位冻结 ⟹
发散轴沿锁定相位的一一映射与压缩（pat
原生观测）}\label{ux516dux6df1ux5165r162ux76f8ux4f4dux51bbux7ed3-ux53d1ux6563ux8f74ux6cbfux9501ux5b9aux76f8ux4f4dux7684ux4e00ux4e00ux6620ux5c04ux4e0eux538bux7f29pat-ux539fux751fux89c2ux6d4b}

\begin{quote}
承接用户指令（2026-08-13）：``既然是相位冻结，就会（让）发散轴沿着被锁定的相位进行一一映射和压缩，观测这个过程，分析可能有多少映射存在，然后观测周期轴上已知的数据和现象情况。''
\end{quote}

\textbf{★核心（pat 原生范式------2026-08-13 命名纪律：无 sqrt、无
i，实数 + Nat + pat 格点）：冻结 = 链方向 d = 0 = pat0 吸收（1
个压缩映射：全链 → 基点）；压缩 = 模 N 周期类（第 k+N
步差一整圈，纤维可数）；一一映射 = 基本区间单射（N 步 ↔ N
槽环，R050）。}

\subsubsection{论证（PatFreezeCompression.lean，7 定理，0
sorry）}\label{ux8bbaux8bc1patfreezecompression.lean7-ux5b9aux74060-sorry}

\begin{enumerate}
\def\labelenumi{\arabic{enumi}.}
\tightlist
\item
  \textbf{冻结 = pat0 吸收}（R137/R122/R134）：冻结（经典 YM
  无质量，相位不演进）= 链方向 d =
  0------每步不移动，整条发散轴（链的全体项）压缩到基点
  pat0（\texttt{freeze\_is\_pat0\_absorbing}：patChain pat0 0 n =
  pat0）。这正是 pat0 吸收（R134：pat0 上每操作 = pat0
  自身）的链表达。\textbf{映射计数：1 个压缩映射}。
\item
  \textbf{压缩 = 模 N 周期类}（R141/R150）：锁定相位 = 2π·j/N（pat
  量化格点）；链第 k 步相位 = 2π·k/N；第 k+N 步与第 k
  步差一整圈（\texttt{period\_class\_full\_turn}：相位差 =
  2π）------压缩纤维 = 模 N 周期类 \{k + m·N\}（可数无限）；槽间距 2π/N
  \textgreater{} 0（\texttt{slot\_spacing\_positive}，一步一个槽）。
\item
  \textbf{一一映射 = 基本区间单射}（R050）：链在 N 步内（j ≠ k
  \textless{}
  N）相位互异（\texttt{slot\_phases\_distinct}）------一个周期（N 步 =
  基本区间）与 N 槽环一一对应；发散轴 = ℵ₀
  个基本区间的并（每区间与槽环一一对应）。
\item
  \textbf{映射计数汇总}：冻结（d = 0）= \textbf{1 个压缩映射}（全链 →
  基点，R134）；锁定（N 槽环）= \textbf{N 个相位槽}（每槽纤维 = 模 N
  可数类）；一一映射存在（基本区间单射，R050）。
\item
  \textbf{周期轴已知现象}（观测清单，全部锚已验收）：3 槽环 ω³ =
  1（PatThreeBody/R141）、临界线圆过 0 和 2（R145）、质量相位周期 T =
  2π/m 往返（R160）------\texttt{period\_axis\_known\_data}。
\end{enumerate}

\textbf{验证}：lake env lean 编译通过，0 error 0 sorry（纯 ℝ + Nat，无
Complex.I）；heap\_verify 枚举 2/2 PASS（周期类代数 / 0
吸收）。★范式纠偏：首版用 exp(2πt/T·i)（分析范式）被用户质疑''这是 pat
的原生研究范式吗''------重写为 pat 原生（R137 链 / R050 单射 / R141 槽环
/ R150 格点 / R134 吸收），证明一次通过，无 cast 调试。

\textbf{诚实边界}：pat 结构侧写（冻结压缩的格点计数），非量子 YM
场论------质量谱存在性（千禧年问题）未触及，R160 边界延续。

\begin{center}\rule{0.5\linewidth}{0.5pt}\end{center}

\subsection{七、相关对照}\label{ux4e03ux76f8ux5173ux5bf9ux7167}

\begin{quote}
完整书目见习题 I
\texttt{../shared/pat\_pnp\_shared\_references\_ZH.md}（共享引用清单）（全字段
+ 验证记录）。本条习题新增对照：
\end{quote}

{\def\LTcaptype{none} % do not increment counter
\begin{longtable}[]{@{}
  >{\raggedright\arraybackslash}p{(\linewidth - 4\tabcolsep) * \real{0.3333}}
  >{\raggedright\arraybackslash}p{(\linewidth - 4\tabcolsep) * \real{0.3333}}
  >{\raggedright\arraybackslash}p{(\linewidth - 4\tabcolsep) * \real{0.3333}}@{}}
\toprule\noalign{}
\begin{minipage}[b]{\linewidth}\raggedright
文献
\end{minipage} & \begin{minipage}[b]{\linewidth}\raggedright
对应内容
\end{minipage} & \begin{minipage}[b]{\linewidth}\raggedright
备注
\end{minipage} \\
\midrule\noalign{}
\endhead
\bottomrule\noalign{}
\endlastfoot
\textbf{Yang, C. N., Mills, R. L.} 1954. \emph{Conservation of Isotopic
Spin and Isotopic Gauge Invariance}. Physical Review 96(1):191--195. &
杨-米尔斯规范理论原始论文 & 无质量规范场（经典尺度不变） \\
\textbf{Jaffe, A., Witten, E.} 2006. \emph{Quantum Yang--Mills Theory}.
In \emph{The Millennium Prize Problems}, Clay Mathematics Institute. &
千禧年问题官方陈述（存在性 + 质量间隙 Δ \textgreater{} 0） &
诚实边界：开放问题，本题只做结构侧写 \\
\textbf{Wilson, K. G.} 1974. \emph{Confinement of quarks}. Physical
Review D 10(8):2445--2459. & 格点规范：质量间隙/色禁闭的数值证据方向 &
证据方向，非数学证明 \\
\textbf{R047/R085/R139/R141/R110}（筑基篇） & 发散-周期特征分解 / 折叠类
0 / 相位-数值互锁 / 单位根相位圆 / log 镜像 & 本框架定理，非外部引理 \\
\end{longtable}
}

\begin{center}\rule{0.5\linewidth}{0.5pt}\end{center}

\emph{筑基篇课后习题 V · 2026-08-13 · Lean 4 / mathlib v4.32.2 · 14
theorems PROVED no sorry（PatYangMills.lean 7 +
PatFreezeCompression.lean 7，R160/R162）· 配套 Lean
形式化：formal/Formal/Toolkit/PatYangMills.lean +
PatFreezeCompression.lean（快照见 ../lean/）}

\end{document}
